% 06_discussion.tex
\section{Conclusions}
\label{sec:discussion}

To conclude, we have formulated QSP in the stereographic
coordinate $x=\cot\theta$, which covers the whole real line, and
have used it to prepare functions on the real line as quantum
states without truncation.  In this coordinate a QSP sequence is
the standard sequence conjugated by a fixed gate, so that its
phases follow from the standard theory.  An unbounded value is
read out as the ratio of two bounded amplitudes, with a number of
shots which grows as the fourth power of that value.  For functions of a block-encoded matrix, the expansion
coefficients and the approximation error at every degree coincide
in every coordinate, so that the known bounds apply unchanged.

We have found that the benefit of the coordinate lies in the
choice of sample points for state preparation.  On the cotangent grid, a function weighted
by the square root of the grid spacing becomes a function on a
uniform grid of angles, and the prepared state represents the
original function in the norm of $L^2(\R)$ with the approximation
error as its only error.  The state is the quantum Fourier
transform of a short coefficient vector and needs no ancilla.
Alternatively, it is prepared by a QSP circuit whose signal is the
Cayley transform of the position operator of the grid.

We have shown that the length of the coefficient vector does not
depend on the error for the Lorentzian and the related amplitudes with algebraic tails, for
which a uniform grid needs a degree that grows as a power of the
inverse error.  For the hyperbolic secant and the Gaussian, in
turn, the length grows polylogarithmically with the inverse error,
and the cotangent grid is more economical for the former and
equivalent for the latter, where its cost with amplitude
amplification lies between $0.9$ and $1.3$ times that of the QSVT
pipeline.

Three properties of the construction set the scope of these
results.  First, the state
represents the function by interpolation in the rational Chebyshev
basis and not by samples on a uniform grid.  A single qubit, for
instance, already represents the Lorentzian exactly.  Hence, a
task which needs the amplitudes on a uniform grid requires a
regridding step, which we leave for future work.
Second, the grid is dense at one point only.  Two Lorentzians of
width $b$ separated by $2|a|\gg b$ therefore need a degree which
grows as $|a|/b$, against $d=1$ for a single one.  Third, the construction requires a signal
operator whose eigenphases we choose.  For a block-encoded
Hamiltonian with given eigenvalues the Cayley transform costs a
linear solve~\cite{DelCorso2026}, and the approximation error
is then the same in every coordinate (Sec.~\ref{sec:limits}).

A natural next step is to identify tasks which consume the state
through the overlaps and expectation values of
Eqs.~\eqref{eq:quadrature} and~\eqref{eq:expectation}, which are exact integrals over the
real line.  One candidate is amplitude estimation of expectation
values under Cauchy and Student $t$ distributions, whose square
roots are the states prepared exactly in
Sec.~\ref{subsec:examples} and for which a uniform grid requires
truncation.
